\section{Introduction}
\label{sec:introduction}

When an agent accepts a supplier's offer, what evidence establishes that its organisation has made an authorised commitment? A conversation may contain the agreed price, a compliance assessment and an instruction to proceed, yet these messages alone do not establish that the assessment covers the accepted terms or that the required resources were reserved. The difficulty becomes more apparent when several organisations operate the agents, retain different records and observe decisions at different times. Agent-to-business (A2B) coordination must connect what agents say to the authority and business state on which their decisions depend.

Consider an airline seeking accommodation after a flight disruption. Its agent discovers hotels, exchanges offers and requests approval within a delegated budget. A hotel agent must allocate inventory before confirming. A revised offer may invalidate an earlier assessment, and a delayed confirmation may leave the airline uncertain whether it is already committed. These are established problems of negotiation and distributed decision-making. Agent mediation brings them together in workflows where participants and delegated tasks may change during the exchange. Earlier business-message standards also show that common formats do not remove the organisational work of connecting trading partners~\citep{iacovou1995edi}.

Current protocols provide useful parts of this workflow. MCP supplies tool and resource access, A2A supports agent interaction, commerce protocols describe commercial exchanges, and AP2 provides delegated-authorisation mechanisms~\citep{mcp2026spec,a2a2026spec,ucpspec,acpspec,ap2auth}. Applications can combine these interfaces with directories, retained histories and transaction controls. The unresolved question is whether making the interaction itself a shared, attributable record improves collaboration across organisational boundaries, compared with assembling equivalent evidence from separately operated services.

We hypothesise that a shared record of participants, proposals, delegations and outcomes can improve independent reconstruction and recovery when agent participation changes or messages are delayed. This hypothesis concerns the representation and availability of evidence. Authority checks, resource reservations and disclosure controls remain necessary in either design; their presence cannot establish an advantage for the record layer. If equally protected point-to-point coordination reconstructs decisions and recovers from interruptions just as effectively at comparable cost, the proposed advantage would not be supported.

We propose \emph{Murmur}, a coordination method for A2B interactions that combines AT Protocol records with A2A task execution. Within this method, organisational gateways bind exact commercial terms to delegated authority and resource reservations, while a release interface controls outbound disclosure. The coordination components are implemented in a sandbox that publishes routing decisions and results and reconstructs conversation context. The commercial and disclosure controls remain proposed. An airline--hotel example makes the proposed transitions explicit. The available implementation and bounded design checks establish starting points for investigation; they do not demonstrate the central comparative hypothesis.

The contribution is a research proposition with an inspectable mechanism and a path to refutation. We first derive the coordination requirements and state the hypothesis, then explain Murmur and the worked transaction. We distinguish the available evidence from the experiments needed to test the proposal, including matched coordination alternatives, interrupted publication and controlled disclosure. The aim is to identify when shared agent conversations become useful evidence for business commitments, and when their operational and privacy costs outweigh that benefit.
